\section{Conclusion}
\label{sec:conclusion}

The story began with a simple bargain: concentrate liquidity in a price range,
earn more fees while the price stays inside it. The bargain hides a cost. The
same position that earns the fees is the position that loses to informed order
flow, and on real pools that loss can exceed the income. An LP who wants to do
better than accept it must decide when to recenter, how wide to stay and how
much risk to hedge, and each decision is expensive and uncertain.

This report set out to measure whether a risk-sensitive learned policy can make
those decisions jointly better than the best rules written by hand, and to find
where active management stops paying. Answering that required a world in which
costs are measured rather than assumed, so we built a testbed from two years of
recorded trades, with an exact replay engine for ground truth and a classical
simulator for training, reconciled by a parity harness. The study then asked
four questions: where fees fail, why a standard agent fails, whether a learned
policy wins out of sample, and where the boundary between active and passive
provision lies.

\pending{This paragraph will close with the measured verdicts once the
experiments complete. Its shape is fixed: what the four answers were, which of
them held, and what the boundary looked like.}

The boundary of the study is also a map for future work. The action space can be
extended with continuous width and an explicit exit decision. The boundary sweep
can cover more pairs and fee tiers. Latency, private order flow and the exact
funding of a real venue can be modelled rather than proxied. Each extension is a
re-run of the same testbed rather than a redesign of it, which is the point of
building the testbed before asking the question.
